\paragraph{Proposition B.6.10 (Proved: the scalar-return cover and its elliptic surface).}
Fix $p\in\mathbb P^1\setminus\{0,1,\infty\}$ and let $C_p$ be
the smooth projective normalization of
\[
 h^{12}=\frac{t^4}{(t-1)^3(t-p)}.
 \tag{B.C1}
\]
Multiplying the right side by a nonzero constant does not affect
the statements below.  Write $\eta(h)=\zeta_{12}h$ for its deck
generator.  This is a connected genus-three cover, with
ramification indices $3,4,12$ at $0,1,p$, respectively.
Let $S_p\to C_p$ be the relatively minimal elliptic surface obtained
by pulling back Engel's marked surface $S$ of type $IV^*,III,I_1$.
Then $S_p$ has precisely twelve singular fibres, all of type $I_1$,
and
\[
 \chi(\mathcal O_{S_p})=1,\quad q(S_p)=p_g(S_p)=3,\quad
 (b_2^+,b_2^-)=(7,15),\quad \sigma(S_p)=-8.
 \tag{B.C2}
\]
The section induced by $P$ has height two, meets the zero section
with intersection number zero, and satisfies $(P-O)^2=-2$.
For $k=0,\ldots,11$, the dimensions of the $\zeta_{12}^k$
eigenspaces of $\eta^*$ on $H^2(S_p;\mathbb C)$ and on the
complexification of its positive subspace are, respectively,
\[
 \begin{array}{c|rrrrrrrrrrrr}
 k&0&1&2&3&4&5&6&7&8&9&10&11\\\hline
 \dim H^2_k&3&1&2&2&2&1&3&1&2&2&2&1\\
 \dim H^+_k&1&0&1&1&0&0&2&0&0&1&1&0 .
 \end{array}
 \tag{B.C3}
\]
In particular the invariant intersection form has rank three
and signature minus one.  This proposition computes the cover
of the scalar return gauge in B.6.3.  It does not identify that
cover with a complete semistable component normalization.

\begin{proof}
The three valuations of the function in (B.C1) are $4,-3,-1$.
Their greatest common divisor with twelve is one, which proves
connectedness.  At these values there are respectively $4,3,1$
points above the base, with the displayed ramification indices.
Riemann--Hurwitz gives
\[
 2g(C_p)-2=-24+(12-4)+(12-3)+(12-1)=4.
\]
There are twelve unramified points over infinity.

The order-three and order-four base changes at $0$ and $1$
are exactly the good-reduction charts used in D.3.3:
$t=z^3,x=z^4X,y=z^6Y$ at zero and
$t-1=z^4,x=z^2X,y=z^3Y$ at one.  Thus the relatively
minimal pullback is smooth over those points.  The fibre at
$p$ was already smooth, so ramified base change there remains
a smooth elliptic fibration.  Over infinity are twelve
unramified copies of $I_1$.  Euler localization gives
$c_2(S_p)=12$.  The canonical bundle of an elliptic surface
with a section and no multiple fibres is pulled back from
its base, so $c_1^2=0$.  Noether's formula yields
$\chi(\mathcal O)=1$.

For completeness the usual direct-image calculation is enough
to determine the irregularity: $\pi_*\mathcal O=\mathcal O_{C_p}$
and $R^1\pi_*\mathcal O=\mathcal L^{-1}$, where
$\deg\mathcal L=\chi(\mathcal O)=1$.  The negative-degree
line bundle has no section, so $q=g(C_p)=3$; hence $p_g=3$.
The Euler number gives $b_2=22$, and the K\"ahler decomposition
gives $b_2^+=1+2p_g=7$.  This proves (B.C2).
Shioda's base-change identity
\cite[Proposition 8.12]{mapsShioda90} multiplies the original
height $1/6$ by twelve.  Since all fibres are irreducible,
his height formula \cite[Theorem 8.6]{mapsShioda90} becomes
$\langle P,P\rangle=2+2(P\cdot O)$; it gives $P\cdot O=0$.
Both section self-intersections are $-\chi=-1$, proving the
last intersection assertion.

We compute the Hodge characters, including their signs.
Choose the rational two-form $\Omega_S$ whose divisor on the
original rational elliptic surface is minus the fibre at infinity.
Its pullback to the minimal surface has orders $0$ at each of the
four points above zero, $2$ at each of the three points above one,
$11$ at the point above $p$, and $-1$ at the twelve points above
infinity.  The first two values follow by substituting the two
good-reduction charts in $dt\wedge dx/y$; the third is the
ramification factor of $dt$, and the last is unramified pullback.
At every other fibre the form is regular and nonvanishing as
a canonical generator.  Its divisor on the base therefore has
degree $6+11-12=5$, as also follows from $K_{C_p}\otimes\mathcal L$.

Every eigenform is a rational function of $t$ times
$h^j\Omega_S$, for some $0\leq j\leq11$.  Its least allowed
factors $t^a(t-1)^b(t-p)^c$ satisfy
\[
 a=-\lfloor j/3\rfloor,\qquad
 b=\lceil(j-2)/4\rceil,\qquad
 c=\lceil(j-11)/12\rceil=0.
\]
Regularity at infinity allows an additional polynomial of
degree at most $-a-b-c-1$.  This degree is nonnegative exactly
for $j=6,9,10$, and is zero in each case.  A basis is consequently
\[
 \frac{h^6(t-1)}{t^2}\Omega_S,\qquad
 \frac{h^9(t-1)^2}{t^3}\Omega_S,\qquad
 \frac{h^{10}(t-1)^2}{t^3}\Omega_S .
 \tag{B.C4}
\]
The invariant K\"ahler class and these forms together with their
conjugates give the positive row of (B.C3).

We determine the full second-cohomology characters by fixed points.
For $1\leq j\leq11$, the number of fixed points of $\eta^j$ on
$C_p$ is
$1+4\mathbf 1_{4\mid j}+3\mathbf 1_{3\mid j}$.
Thus its first-cohomology trace is two minus this number.
Pullback identifies $H^1(C_p)$ with $H^1(S_p)$: it is injective
by the zero section and both groups have dimension six.
The traces on $H^1$ and $H^3$ agree by duality and conjugation.
Above $p$ the whole smooth elliptic fibre is fixed and has Euler
number zero.  At a point above zero, a nontrivial order-three
elliptic automorphism fixes three points.  At a point above one,
an order-four automorphism fixes two points, and its square fixes
four.  There are no other fixed points.  The local Lefschetz signs
are positive, since the normal actions are complex rotations.
Subtracting the contributions of $H^0,H^1,H^3,H^4$ gives
\[
 \operatorname{Tr}(\eta^{j*}\mid H^2)=
 \begin{cases}
 22&j=0,\\
 4&j=4,8,\\
 6&j=6,\\
 0&\text{otherwise}.
 \end{cases}
\]
The finite Fourier transform gives the first row of (B.C3).

The program \texttt{fourfolds\_B6\_cover.py} performs these twelve
character calculations in $\mathbb Q[x]/(x^4-x^2+1)$ and checks
every displayed dimension exactly; no numerical eigenvalues are used.
\end{proof}

\paragraph{Proposition B.6.11 (Proved: the untwisted plumbing has signature two).}
Let $W^\pm$ be a sufficiently small closed parameter-disc
neighbourhood in either smooth semistable total fourfold of
K13, A.2.2, with its complex orientation.  Its boundary is the
mapping torus of the \emph{untwisted} family $X(u)$ under one
positive $u$ loop, and this boundary is not spin.  One has
\[
 \sigma(W^\pm)=2,\qquad c_1(W^\pm)^2=0
       \quad\text{in }H^4(W^\pm;\mathbb Z).
 \tag{B.P1}
\]
The signature form modulo its radical is rationally represented by
$\left(\begin{smallmatrix}2&1\\1&6\end{smallmatrix}\right)$.
In particular the $E_8$ summand of the rational elliptic surface
does not give an $E_8$ signature for this parameter filling.
These are not the finite logarithmic replacements used for the
sphere; no smooth monodromy value for those replacements follows.

\begin{proof}
We first give the local-to-global intersection calculation.
Write $Q$ for the irreducible central normal-crossings fibre and
$\nu:R\to Q\subset W$ for its smooth normalization; here $R$ means
$R^+$ or $R^-$.  Its conductor sections $i_0,i_\infty:S\to R$
are glued by $T=t_{6P}$.  A sufficiently small semistable tube
retracts to its central fibre: in the product charts $xy=u$ shrink
the normal coordinates toward their zero loci, fixing that locus,
and use positive radial normal fields on the single-branch charts.
Collars make these contractions agree, and compactness gives a
finite choice of charts.  The normalization mapping cylinder
therefore computes the homology of $W$.  In particular
\[
 H_4(S;\mathbb Q)\longrightarrow H_4(R;\mathbb Q)
 \longrightarrow H_4(W;\mathbb Q)\longrightarrow H_3(S;\mathbb Q)=0.
 \tag{B.P2}
\]
Thus all middle classes are pushed forward from the normalization.

The normalization is a complex codimension-one immersion with
transverse double locus.  Let $D_0,D_\infty$ be its two conductor
divisors.  Its normal line has
\[
 c_1(N_\nu)=-[D_0]-[D_\infty].
\]
Indeed the parameter differential in $u=xy$, restricted to either
branch, is the normal differential multiplied by the coordinate
vanishing on the conductor.  Since the parameter divisor is
principal, these local identities identify the normal line with
$\mathcal O_R(-D_0-D_\infty)$.

For $a,b\in H^2(R;\mathbb Q)$, put $a_i=i_i^*a$ and $b_i=i_i^*b$.
Push the four-cycles Poincar\'e dual to $a,b$ into $W$, and perturb
one transversely.  Their same-branch intersections are counted by
the Euler class of $N_\nu$; their two other intersections lie on
the two branches of the conductor.  The complex orientations give
positive signs to both cross terms.  Thus their intersection is
\[
 -a_0b_0-a_\infty b_\infty
       +a_0T^*b_\infty+b_0T^*a_\infty
   =-\bigl(a_0-T^*a_\infty\bigr)
         \bigl(b_0-T^*b_\infty\bigr),
 \tag{B.P3}
\]
where products on the right are evaluated on $[S]$.
This also follows by doing the perturbation in each $xy=0$ chart;
the two unequal branch intersections are the two last terms.
Define $U(a)=a_0-T^*a_\infty$.
Equation (B.P2) shows that the signature is precisely the negative
of the surface intersection form restricted to $\operatorname{im}U$.
The relation coming from $H_4(S)$ causes no discrepancy:
its dual divisor is $D_0-D_\infty$, and its image under $U$ is zero
because the two conductor normal bundles have trivial product.

We compute this image using the actual surface classes
\[
 F=[S_t],\qquad m=[P]-[O],\qquad
 v=[6P]-[O]-3F.
\]
They satisfy
\[
 F^2=Fm=Fv=0,
 \qquad m^2=-2,\quad mv=-1,\quad v^2=-6.
 \tag{B.P4}
\]
These equalities follow from the section intersections
$P\cdot O=0$, $[6P]\cdot O=P\cdot[6P]=2$ and the
self-intersection $-1$ of every section.  The two latter
intersection numbers follow directly from height $1/6$ and
the narrow section $6P$, as in \cite[Lemma 2.1]{mapsEngel}.
Translation acts by the integral isometry
\[
 T_*x=x+(xF)v-(xv-3xF)F.
\]
It is checked on $F,O,v$ and the seven component-root classes,
which $t_{6P}$ fixes.  In particular
$\operatorname{im}(I-T^*)_{\mathbb Q}=\mathbb QF+\mathbb Qv$.
On $\mathbb P(\mathcal O\oplus L)$, the relative hyperplane
class has two section restrictions differing by $m$; hence its
image under $U$ is $m$ modulo $\mathbb QF+\mathbb Qv$.
The blow-up in K13 is along the fibre $S_\infty$ in one section.
Its exceptional divisor restricts to $F$ on that section and zero
on the other.  The small resolution introduces no further divisor
class: the normalization homology, or the two displayed conifold
charts, gives one additional second-cohomology generator in total.
Consequently
\[
 \operatorname{im}U=\mathbb QF+\mathbb Qv+\mathbb Qm.
\]
Its radical is $\mathbb QF$.  The matrix on $(m,v)$ is negative
definite, with leading entry $-2$ and determinant eleven.
Negating it in (B.P3) gives the positive matrix in the statement,
and proves $\sigma(W)=2$.  The signs and determinant are checked
exactly in \texttt{fourfolds\_B6\_plumbing.py}.

Finally on the initial projective bundle,
$c_1(TR)=\pi_R^*c_1(S)+D_0+D_\infty$.
Blowing up the conductor curve subtracts the exceptional divisor
both from $c_1(TR)$ and from $D_0+D_\infty$; the subsequent small
conifold resolution is crepant.  The tangent-normal sequence
therefore gives
\[
 \nu^*c_1(W)=\pi_R^*c_1(S)=\pi_R^*F.
\]
Its square is zero.  The cohomological normalization sequence,
with $H^3(S;\mathbb Z)=0$, makes
$H^4(W;\mathbb Z)\to H^4(R;\mathbb Z)$ injective.
This proves the integral assertion in (B.P1).
We check the boundary assertion integrally as well.
The normalization sequence gives $H_6(W;\mathbb Z)=\mathbb Z$,
generated by the reduced irreducible central fibre, and
$H_5(W;\mathbb Z)=0$.  For the latter, $H_5(R)=0$ and the next
map $H_4(S)\to H_4(R)$ is nonzero: after pushing through the
normalization modification its dual is $m\ne0$.
Poincar\'e--Lefschetz duality identifies
$H^2(W,\partial W)\to H^2(W)$ with the class of that central
divisor.  It is zero because the divisor is the zero set of the
global parameter function.  The next relative group is
$H^3(W,\partial W)=H_5(W)=0$.  The exact sequence thus gives
\[
 H^2(W;\mathbb Z)\xrightarrow{\ \simeq\ }
 H^2(\partial W;\mathbb Z).
\]
The class $c_1(W)$ is primitive: its pullback $\pi_R^*F$ pairs
to one with the zero-section curve in either conductor copy of
$S$.  Its restriction to the boundary is therefore not divisible
by two.  For a complex tangent bundle $w_2=c_1\bmod2$, and the
stable tangent restriction to the boundary has the same $w_2$.
Consequently $w_2(\partial W)\ne0$, proving that this is not the
spin mapping torus required in B.6.9.

When relative characteristic numbers for a compatible boundary
are defined, the absolute vanishing also kills the terms
$c_1^4$ and $p_1c_1^2$ by pairing a relative lift with the zero
absolute square.  It does not supply the missing relative
$p_1^2$ or replace the finite logarithmic caps.
\end{proof}
